\newpage
%%%%%%%%%%%%%%%%%%%%%%%%%%%%%%%%%%%%%%%%%%%%%%%%%%%%%%%%%%%%%%%%
%%%%%%%%%%%%%%%%%%%%%%%%%%%%%%%%%%%%%%%%%%%%%%%%%%%%%%%%%%%%%%%%
%%%%%%%%%%%%%%%%%%%%%%%%%% Enunciado %%%%%%%%%%%%%%%%%%%%%%%%%%%

\begin{myblock}
\phantomsection\addcontentsline{toc}{section}{Ejercicio \#5 | Clasificador de Géneros Musicales}
\section*{Ejercicio \#5 | Clasificador de Géneros Musicales}

En el moodle del curso, encontrarás un conjunto de datos que corresponden a un extracto del Free Music Archive (FMA) [1], que es una base de datos muy extensa de archivos de audio usada para diversas tareas de MIR.\footnote{El Free Music Archive es un repositorio interactivo de música legal y de alta calidad.}

Encontrarás varios conjuntos de datos en formato \texttt{.csv} que corresponden a dos grupos: uno de entrenamiento (archivos que empiezan con el prefijo \texttt{tracks\_fma\_train}) y uno de prueba (archivos que empiezan con el prefijo \texttt{tracks\_fma\_test}). Para cada uno hay un archivo con metadatos (\texttt{\dots\_metadata.csv}) que incluye la variable de respuesta \texttt{track.genre12} excepto para los datos de test.

También, se proporcionan diversos archivos que contienen características de audio y de la señal (\texttt{\dots\_features.csv}) en forma de indicadores o estadísticas de la señal, por lo que las escalas pueden variar.\footnote{Las características se extraen utilizando herramientas estándar de análisis de audio.}

\begin{itemize}
    \item Usando los datos de \textit{train}, realiza un análisis exploratorio de los datos, usando métodos de reducción de dimensión que consideres apropiados. ¿Qué características tiene la variable de respuesta? ¿Puedes identificar los géneros, o al menos, algún subconjunto de ellos? Utiliza las características de audio y de la señal que creas convenientes y repórtalo junto con todos tus hallazgos.

    \item Usando también los datos de entrenamiento, ajusta una red neuronal con \texttt{pytorch} para predecir el género. Genera un conjunto de datos de entrenamiento y validación para verificar su ajuste. Define la arquitectura de la red, y reporta todos los detalles del modelo y el preprocesamiento de los datos que hayas usado (características de audio y señal usadas, reducción de dimensión, regularización, optimizador, estandarización, escalamiento, etcétera). Reporta las métricas de desempeño y gráficas que creas conveniente.

    \item Usando la red neuronal ajustada en el paso previo, predice el género de los datos de \textit{test}.
\end{itemize}

\end{myblock}

%%%%%%%%%%%%%%%%%%%%%%%%%%%%%%%%%%%%%%%%%%%%%%%%%%%%%%%%%%%%%%%%
%%%%%%%%%%%%%%%%%%%%%%%%%%%%%%%%%%%%%%%%%%%%%%%%%%%%%%%%%%%%%%%%

%%%%%%%%%%%%%%%%%%%%%%%%%%%%%%%%%%%%%%%%%%%%%%%%%%%%%%%%%%%%%%%%
%%%%%%%%%%%%%%%%%%%%%%%%%%%%%%%%%%%%%%%%%%%%%%%%%%%%%%%%%%%%%%%%

Los resultados de este problema se encuentran en los notebooks:

\begin{itemize}
    \item \textbf{MLP-MIR-Smooth.ipynb}
    \item \textbf{exploratorio-mir-tarea01.ipynb}
\end{itemize}

\clearpage